\subsubsection{24.03.15}

Today we had a one-day trip to Moscow. In Moscow we had been at the event "PTC Live Tech Forum" organised by our sponsors, company PTC. We were invited there by PTC to show how we use their software. There we presented our model of robot assembled in Creo Parametric 3.0.\newline
The second reason to visit Moscow was a number of special TETRIX detailes, which we ordered before. To take them we needed to meet the representative of PTC, Dmitriy Orlov.\newline
Also the forum was a good chance for us to meet new people who could cooperate with us or become our sponsors.\newline
During the forum, we continued working on our robot (we took it to moscow with us).

\begin{enumerate}
	
	\item The time of beginning and ending of the meeting: 8:00 - 18:20.
	
	\item Purposes of the meeting: 
	\begin{enumerate}
		
		\item Install to the wheel base AndyMark motors.
		
		\item Hold the wiring.
		
        \item Install new longer beams to the mount for controllers.
        
        \item Stregthen the wheel base with new plates.
        
        \item Change the program and test robot's motion.
		
	\end{enumerate}

	\item Work that has been done:
	\begin{enumerate}
		
		\item We installed to the wheel base Andymark motors (which were in our order) instead of DC. Next we connected all the motors and encoders to the Hi-technic controllers.
		
		\item Unfortunately, in the order were no special beams which we wanted to install on wheelbase to make it more reliable.
		\begin{figure}[H]
			\begin{minipage}[h]{0.2\linewidth}
				\center  
			\end{minipage}
			\begin{minipage}[h]{0.6\linewidth}
				\center{\includegraphics[scale=0.3]{days/24.03.15/images/01}}
				\caption{}
			\end{minipage}
		\end{figure}
		
        \item However, we got other long beams which we needed. So we installed them on the mount for controllers. Since then the construction became more straight because there is only one long beam and not the several short ones.
        \begin{figure}[H]
	  	  \begin{minipage}[h]{0.2\linewidth}
	  	    \center  
	  	  \end{minipage}
	  	  \begin{minipage}[h]{0.6\linewidth}
	  		\center{\includegraphics[scale=0.3]{days/24.03.15/images/02}}
	  		\caption{}
	  	  \end{minipage}
	   \end{figure}
	   
	   \item After we finished working on construction, we decided to make a trial run of the robot. To do this, we changed our program a bit: inversed the direction of rotating of two motors (one from each side), because the motors linked with gears and need to rotate in opposite directions. During the testing we noticed, that robot's movement became easier to control: as the center of gravity is now at the same distance from each wheel, it turns around one point. As a result, it moves more predictable.
	   
	   \item We decided to remake the mechanism of overturning of the bucket because of weakness and unreliability of present construction. Firstly, we discussed the general aspects of new construction. The material will be the aluminium L-shaped profile 5x5 cm. After the conversation we created an assembly of new module in Creo Parametric.
	   \begin{figure}[H]
	   	\begin{minipage}[h]{0.2\linewidth}
	   		\center  
	   	\end{minipage}
	   	\begin{minipage}[h]{0.6\linewidth}
	   		\center{\includegraphics[scale=0.3]{days/24.03.15/images/03}}
	   		\caption{}
	   	\end{minipage}
	   \end{figure}

	\end{enumerate}
	
	\item Results:
	\begin{enumerate}
		
		\item AndyMark motors are installed to the wheel base.
		
		\item All the wiring is held.
		
		\item New longer beams were installed to the mount for controllers.
		
		\item The flats weren't fixed on the wheel base.
		
		\item We tested how the robot moves. Results are positive.
		
		\item We created a model of new MOB.
		
	\end{enumerate}
	
	\item Tasks for the next meetings:
	\begin{enumerate}
		
		\item Make new MOB.
		
		\item Make blades for gripper for balls and the ramp for guiding balls into the bucket.
		
		\item Finish working on plexiglass protection.
			
	\end{enumerate}
\end{enumerate}
\fillpage
